\label{chap:conclusions}

\section{Answers to the research questions}
\noindent\textbf{RQ1: Why do small VLMs fail as navigation agents?} Because the agents they are placed in make them
keep spatial state in generated text. Controlled experiments on the same episodes ruled out the two other causes: a
geometric planner given the target position solves \nOracle\% of the EB-Navigation test episodes, and Qwen3-VL-4B
places a correct box on a visible named target in \nProbeThreeCommon--\nProbeThreeBase\% of single views. The
failures of the prompted agent are failures of state: it loses seen targets, collides with furniture, and does not
turn toward a target behind it (\cref{chap:diagnosis}).

\noindent\textbf{RQ2: What should an agent built around a small VLM look like?} It should ask the model only what
the current image can settle and keep everything else in metric form. \ours does so with two per-image questions, a
hypothesis memory, a search rule, and an occupancy-map planner (\cref{chap:method}). Of its components, multi-view
evidence fusion carries the gain: a generic modular agent that acts on single-view answers trails \ours by
\nSOverMod{} percentage points, and removing fusion alone costs \nAblFusionAbs{} (\cref{sec:exp-modular,sec:exp-ablation}).

\noindent\textbf{RQ3: How far does this agent close the gap to larger models?} With a 4-bit 4B model on one 8\,GB GPU
and without training, \ours succeeds in \nSThree\% of the EB-Navigation test episodes, \nGainThree{} percentage
points more than the benchmark's agent with the same model. From RGB alone it reaches \nRGBThree\%, which is above
both fine-tuned 4B agents and seven of ten training-free agents built on 8B to frontier-scale models. With a smaller
3B model the gain grows to \nGainTwo{} percentage points, and on RoboTHOR \ours reaches \nRoboSR\% success, above ten
published zero-shot methods (\cref{chap:experiments}).

\section{Discussion}
The central finding of the thesis is that, for small models, the role matters more than the model. Agents designed
around frontier models ask the model to perceive, remember, and plan, and they degrade quickly as the model shrinks.
An agent that asks the same model only what it can see loses far less: replacing the 4B model by a 3B model costs
\ours \nScaleDrop{} percentage points and the benchmark's agent \nScaleDropEB. The division of labor between a
semantic sensor and a metric controller is old in robotics, but the thesis shows that a generative VLM needs more
than this division: because its single answers are unreliable, the controller must treat them as evidence and
accumulate them over distinct viewpoints.

\subsection{Implications for the design of robot agents}
Three lessons from this thesis apply beyond EB-Navigation.

\noindent\textbf{Match the question to the model.} A question about the current image is one a small model can
answer and an agent can check. A question about the room, the history, or the plan is one it cannot. Agent designs
that pass the whole problem to the model test the model's ability to plan, not its ability to see. For small models
the second is far stronger than the first.

\noindent\textbf{Treat a generative model as a noisy sensor.} A generative VLM places boxes in views that do not
show the target and repeats its answers from the same viewpoint. An agent should therefore weigh its answers, combine
them across viewpoints, and lower them when the expected detection does not appear, as one would with any other
noisy sensor. Of all components of \ours, this evidence fusion is the one whose removal costs most.

\noindent\textbf{Keep geometry in geometry.} Distances, free space, and directions already looked at are exact
quantities that classical methods compute reliably from depth and pose, or from RGB alone with a ground-plane scale.
There is no reason to ask a language model to estimate them in words.

\section{Limitations}
\noindent\textbf{Task scope.} The controller of \ours is written for one task, reaching one object. A prompted agent
can be redirected to a new task by changing its prompt; \ours needs a new controller.

\noindent\textbf{Scene scale.} EB-Navigation episodes take place in single rooms with at most 20 steps. On the
multi-room RoboTHOR apartments \ours remains below two zero-shot methods (\cref{tab:robothor}). Its exploration
follows geometric rules and does not use semantic priors about where objects tend to be, which matter more in larger
scenes. The components designed for competing candidates showed no significant effect in single rooms and remain
untested at a larger scale.

\noindent\textbf{Recognition.} Visual is the only subset on which \ours does not exceed the benchmark's agent. When the
instruction describes only appearance, the small model rarely proposes the target, and no memory can recover a target
that is never detected (\cref{sec:exp-gap}).

\noindent\textbf{Simulation only.} All results come from AI2-THOR. The RGB-only configuration removes the depth sensor
and the simulator pose, but its dead reckoning benefits from the simulator's discrete actions; real actuation noise,
motion blur, and lighting have not been tested.

\section{Future work}
\noindent\textbf{Deployment on a robot.} The agent needs only an RGB camera, odometry, and a 4B model on a consumer
GPU, which makes on-board deployment on a mobile or legged robot the most direct next step.

\noindent\textbf{Semantic exploration.} The VLM could also answer questions that the current image can settle about
where to search, for example which visible doorway leads toward a kitchen, without being asked to plan.

\noindent\textbf{Larger scenes and other tasks.} Multi-room benchmarks would test the memory and the search index at
a larger scale, and controllers for multi-step tasks, such as finding an object and then a receptacle, would test how
far the division of labor between a semantic sensor and a metric controller extends.

\noindent\textbf{Appearance-only goals.} Comparing several candidates within one image, rather than judging each
alone, may help a small model recognize objects described only by how they look.
